% A dependency audit and two explicit conditional joining lemmas.
% This fragment uses only the macros and theorem environments of the main note.
\section{A dependency audit organized by complete chains}
\label{sec:dependency-audit}

\status{This is a human mathematical audit of the reconstruction and a
targeted reading of the pinned source interfaces.  It is not an extraction
of Lean proof-term dependencies.  Its two joining arguments below are
complete implications from the stated premises.  The arithmetic suppliers
of those premises remain separate obligations.}

There are three different meanings of ``complete'' that matter here.
A local theorem can have a complete proof for every datum satisfying its
hypotheses.  A chain can be complete relative to explicitly named earlier
theorems.  A global application is complete only when all those hypotheses
hold on one actual family, with the required order of choices.  The local
Type-I and Type-II results are examples of the first meaning; QP's
\texttt{PairEx}, relative to CN/BD/HD, is an example of the second.
The remaining work toward the arithmetic window is predominantly the third.

The historical claim graph was designed to record mathematical references,
not to express every retained premise as a graph vertex.  For example,
AN2 has no predecessor vertex but assumes prime windows; AN3 and SC2 have
empty \texttt{remaining\_inputs} arrays although their prose retains
substantive hypotheses.  QP25 records QP22 as a predecessor and also
assumes \texttt{VWinAt}.  None of these is a mathematical contradiction.
It means that graph acyclicity, or the number of recorded claims, cannot
certify completion.  The separate audit data makes such obligations
first-class objects and leaves all historical edges unchanged.

\subsection{The exact top-level sufficient interfaces}

Fix a primitive target $\psi$ and let $H_\psi$ be its fixed positive
character family.  The downstream chain is
\[
 \begin{gathered}
 \text{one compatible family on arbitrarily late windows}\\
 +\ \text{actual arithmetic gap and numerical bounds}\\
 \Longrightarrow H_\psi(x)\le x^{1/2+A_\psi\epsilon}
 \Longrightarrow \text{zero-freeness for the target}.
 \end{gathered}                                                    \tag{DA.1}
\]
Here $A_\psi$ and a positive permitted range for $\epsilon$ are fixed
\emph{before} $\epsilon$ is chosen.  Taking every primitive target and
using GH6 gives the Dirichlet conclusions; taking only the trivial target
gives ordinary RH.  The full positive-family analytic implication is
AN2--AN3, not an additional open analytic problem.

The literal source organizes the upstream cut using
\texttt{GChi}$(\psi)$, which requires, on the same host and selected set,
\[
 \texttt{EAtH}(\operatorname{Level.cyc}\psi,H,T_T,x,\epsilon,4)
 \quad\text{and}\quad
 \texttt{VVle}(E,t,\sqrt b).                               \tag{DA.2}
\]
Its family class also requires arbitrarily late windows.  The first
condition supplies the character estimates and the second the
cohomological budget.  The source's
\texttt{HostExW} supplies the actual ranks, \texttt{HostGood} data and
this target requirement on a scale family.  A description of the abstract
field $E$, or the existence of a pair $q,d$, does not yet provide that
whole record.

The sufficient source cut is
\[
 \begin{gathered}
 \texttt{HostExW}(3/2,\texttt{GChi}(\psi).\texttt{req})\quad\text{and}\quad
 \text{the design/link construction}\\
 \Longrightarrow
 \texttt{FamilyClassConnected}(\texttt{GChi}(\psi).\texttt{cls}).
 \end{gathered}
                                                               \tag{DA.3}
\]
The last predicate supplies fixed $e_D,e_L$ before $\epsilon$ and a
family with \texttt{ChainTied} at every sufficiently late scale in its
window set.  The final character/scalar chain then gives (DA.1).
This is a sufficient cut through the source argument, not a claim of
logical necessity or a minimal set of possible mathematical ideas.

\subsection{How one family can satisfy finitely many closed conditions}

The source has an explicit solution to the order-of-choices problem.
Write $\mathcal C$ for a class of families and
$\operatorname{Large}(F,K,H)$ for its monotone threshold condition.
Thresholds have a finite supremum satisfying
\[
 \operatorname{Large}(F,K\vee K',H)
 \quad\Longleftrightarrow\quad
 \operatorname{Large}(F,K,H)\ \text{and}\
 \operatorname{Large}(F,K',H).                             \tag{DA.4}
\]
The literal definition of \texttt{Closed} is
\[
 \begin{split}
 \operatorname{Closed}_{\mathcal C}(P):\quad
 &\exists\epsilon_P>0\ \forall\epsilon\in(0,\epsilon_P]
 \ \exists K_P\ \forall F,H,\\
 &F\in\mathcal C,\ \operatorname{Large}(F,K_P,H)
 \ \Longrightarrow P(F,H).
 \end{split}                                                   \tag{DA.5}
\]
In contrast, \texttt{DesignExists} quantifies over the threshold first:
for each sufficiently small $\epsilon$ and \emph{every} $K$, it supplies
one $F,H$ in the required class, large for $K$, with
\texttt{ChosenE}, \texttt{TuSize} and \texttt{DGamSize}.

\begin{proposition}[DA1: finite uniform joining on one family]
\label{prop:da-closed-join}
If $P_1,\ldots,P_r$ satisfy (DA.5) on the same family class, then
their conjunction does too.  If \texttt{DesignExists} supplies families
in that class, one chosen family satisfies all $P_i$ and all the design
conditions simultaneously.  The positive range for $\epsilon$ is chosen
before $\epsilon$, and the joined threshold is chosen before the family.
\end{proposition}

\begin{proof}
Include the design supplier's range in the finite minimum of all
positive ranges.  For an $\epsilon$ in that range, take
the threshold $K_i$ furnished by each closed statement and put
$K=\bigvee_iK_i$.  Design existence, which holds for every $K$, now
chooses $F,H$ large for this $K$.  Repeated use of (DA.4) makes it large
for each $K_i$, so every $P_i(F,H)$ follows.  For no conditions there is
nothing to join; choose any design threshold.  All conclusions use the
one family just selected.  This is the entire mathematical content of
the source's \texttt{Closed.and} and its use in \texttt{topDesWin}.
It does not prove \texttt{DesignExists} or any of the closed premises.
\end{proof}

This quantifier structure is stronger than obtaining an individual
witness for each desired property.  It is also more precise than
invoking compatibility after all witnesses have already been chosen.
Every closure must refer to the same class; equal dimensions, isomorphic
fields, or unrelated enlargements of coefficient fields do not supply
the dependent identifications in (DA.5).

\subsection{The complete logical passage through the global links}

Here is the finite conjunction used by \texttt{topDesWin}.  In the table,
``closed'' means the exact quantifiers (DA.5), on the same design class.
Conditional clauses retain their dependence on the design's
\texttt{ChosenE}, \texttt{TuSize} and \texttt{DGamSize} fields.

\begin{center}\small
\begin{tabular}{p{0.24\linewidth}p{0.67\linewidth}}
\toprule
Input & Required content on the actual family\\
\midrule
\texttt{DesignExists} & One target-compatible family after every
finite threshold choice, with the selected parameters and size fields.\\
\texttt{LinkN} & Congruence-denominator cost bounded by the actual
cohomological sum and polynomial budget.\\
\texttt{LinkEta} & Coverage of the actual twisted local data.\\
\texttt{LinkIRR} & Nonnegative retained-row defect and its stated
exponentially small upper bound.\\
\texttt{SmallLevelRankEq} & The actual equality $d_{Z,0}(j)=D_0$
for every lift used.\\
\texttt{LinkDIVt} & Given \texttt{LinkEta} and \texttt{DIVinN}, the
common period target, integral/height bounds and \texttt{LinkDIV}.\\
Fullness and rank & \texttt{ChosenE} implies \texttt{NuFullF} and
positive retained ranks \texttt{J.MuRank}.\\
\texttt{LinkG} & The local divisibility/gain theorem on these actual
lattices and served orbits.\\
\texttt{LinkS} & The actual normalized off-diagonal and diagonal Gram
bounds and total-target dimension bound.\\
\texttt{LinkGLOBAL} & The same-family slope comparison turns
\texttt{ChosenE} and the five global links into \texttt{ChainTied}.\\
\bottomrule
\end{tabular}
\end{center}

The target class supplies \texttt{HypV}, namely
$\texttt{VVle}(E,t,\sqrt b)$ at every sufficiently late scale.
The exact metric transport used by the source is additionally visible:
\[
 \texttt{LinkDIV}+\texttt{NuFullF}+\texttt{J.MuRank}
 +\texttt{J.MuTu}+\texttt{J.MuDGam}
 \Longrightarrow\texttt{LinkMu}.                          \tag{DA.6}
\]
Its final wrapper is \texttt{J.linkMu\_of\_linkDIV}; it relies on the
actual \texttt{NprimeMuR} lattice/metric comparison.  That comparison
is retained here.  In particular, (DA.6) is not being deduced merely
from the ranks of unrelated spaces.

\begin{proposition}[DA2: the conditional link chain closes]
\label{prop:da-topdeswin}
Assume the exact common-class design and closed premises in the table,
the target-class inclusion and \texttt{HypV}, the size consequences
\texttt{TuSize}$\Rightarrow$\texttt{J.MuTu} and
\texttt{DGamSize}$\Rightarrow$\texttt{J.MuDGam}, and the actual
metric transport (DA.6).  Then they supply one family with
\texttt{FamilyClassConnected}, with the constants $e_D,e_L$ in
\texttt{LinkGLOBAL} fixed before $\epsilon$.
\end{proposition}

\begin{proof}
Apply DA1 to the nine closed link/rank conditions in the table and to
the design supplier, after transporting along the stated class-inclusion
maps.  All statements now concern one $F,H$.  It remains to check the
one numerical arrow that turns \texttt{HypV} into the needed
congruence budget; the other arrows are the explicit hypotheses.

At a sufficiently late scale, write $V$ for the actual weighted sum of
the finite cohomological lengths.  The finite-subsums clause in
\texttt{VVle}, together with nonnegativity, shows that this series is
summable and
\[
                     0\le V\le\sqrt b\le b\qquad(b\ge1).
                                                               \tag{DA.7}
\]
The finiteness clause and the literal \texttt{NData.sha\_eq} identify
these lengths with the stored lengths in the same $D$.  Thus
\texttt{ShaFinOdd} holds and \texttt{LinkN} applies.
Put $y=x^\epsilon$, so $y\ge1$, and let
$B_0=fmb\,y$, with $fm\ge1$ and $N_c\ge0$.  Eventually
$\log^2x\le y$.  The \texttt{LinkN} inequality, with
$d_{Z,0}(j)=D_0$ from \texttt{SmallLevelRankEq}, is
\[
 \eta_j\le D_0\{C(fmb\,y^2+\log^2x\,N_c)+V\}.
                                                               \tag{DA.8}
\]
The expression multiplying $C$ is nonnegative and is at most
$y(B_0+N_c)$.  Also $V\le b\le B_0\le y(B_0+N_c)$.
Replacing $C$ by $|C|$ proves
\[
           \eta_j\le D_0(|C|+1)y(B_0+N_c),                \tag{DA.9}
\]
which is precisely \texttt{DIVinN}.  There are only finitely many
eventual statements in this argument, each already uniform in all
family indices and used lifts.  Taking the maximum of their thresholds
therefore preserves that uniformity; it does not assume that the index
sets themselves are finite.

Now \texttt{LinkEta} and \texttt{LinkDIVt} give \texttt{LinkDIV}.
The design size fields give the two size consequences in (DA.6);
fullness and positive rank are supplied on this very family, so (DA.6)
gives \texttt{LinkMu}.  The actual package
\[
 \texttt{GlobalTheoremInputs}
 =\texttt{LinkG}\wedge\texttt{LinkS}\wedge\texttt{LinkIRR}
     \wedge\texttt{LinkDIV}\wedge\texttt{LinkMu}
\]
is therefore available.  Apply \texttt{LinkGLOBAL} with
\texttt{ChosenE}.  This supplies the remaining constants and
\texttt{ChainTied} at every sufficiently late scale in the same
family's window set.  The definition of
\texttt{FamilyClassConnected} is exactly this conclusion, with
$e_D,e_L$ already fixed.  This proves the asserted conditional chain.
\end{proof}

DA2 proves a complete assembly; it does not turn any listed closed
supplier into an existence theorem.  The source's particular
\texttt{shapeWin} call uses $C_{\rm abs}=2,C_0=3,e_D=2,e_L=0$.
Thus its literal final coefficient is $e_D+e_L+7=9$ if that entire
route and all its hypotheses are supplied.  The reconstruction need
not fix this numerical value: any finite coefficient fixed before
$\epsilon$ suffices for the analytic contradiction.

\subsection{The remaining cut and the alternative routes}

The audit uses an AND gate for simultaneous requirements and an OR gate
only when either complete route suffices.  The following are the main
branches; none can be retired solely because its downstream assembly
has been written.

The strongest new local leaf is HP1--HP4: the characteristic-zero
canonical polynomial highest-weight vector, fixed line, and tensor
irreducibility theorem are now proved.  HP5--HP6 identify these with
the existing weight/form interfaces while retaining the exact supplied
field, frame and row data.  Thus the Schur-theoretic premise formerly
left in MR5/MR19 can be discharged.  The geometric realization of that
representation and the other fields of the actual \texttt{Core} are
different obligations.  WB1--WB8 and VW1--VW4 close two complete
downstream interfaces, as specified below; their arithmetic and
cohomological suppliers remain explicit.

\begin{enumerate}
\item \textbf{Cohomological windows.} The complete downstream route is
\texttt{JointWin}+\texttt{ThmBc}$\Rightarrow$\texttt{VWinAt}, then QP22
supplies the fixed-pair windows.  An upstream sufficient alternative
is \texttt{ThmBab}, \texttt{PointHV}, and \texttt{ThmBcFree}.
The actual two-primary and odd-plus/odd-minus bounds remain the leaves.
Separate bounds on separate intervals cannot replace the joint-window
premise.
\item \textbf{Actual geometric and arithmetic family.} Full
\texttt{Core}, the dependent layout, original coefficient embeddings,
\texttt{St}/\texttt{ColPin}/\texttt{R3.Read}, rational good-row input
and genuine theta/local installation remain.  MR22's initial zero
theta and empty toric family are exactly its stated initialization;
they do not supply a later nonplaceholder \texttt{NTh} or local family.
\item \textbf{Actual theta and global analytic operators.} Critical
toric ratios and \texttt{Choice.FieldOK}; convergence, separation,
nonvanishing and operator identities; actual chain norms, normalized
spectral words, positive kernel/mass identities and exceptional terms
are distinct obligations.  The finite Weil dictionary does not imply
these identities for the global operators.
\item \textbf{Integral/metric and congruence comparison.} The true
source/target lattices and metrics, discarded directions, retained
ranks and \texttt{LinkN} denominator bounds must remain on the same
objects.  Rank identities and abstract isometries on different
spaces do not prove the required transport.
\end{enumerate}

There is a useful OR gate within the real-stage branch.  For the actual
large-gap rows, DC20 supplies the four \texttt{RealTied} fields using
the named Harish--Chandra degree theorem and GV's Tamagawa input with
the exact normalizations.  The stronger universal route uses
\texttt{DiscreteSeriesExist} and \texttt{FormalDegreeFormula2}.
It is an alternative and need not also be reconstructed if every row used
is covered by the first route.  In either case the other fields of
\texttt{Core} remain required.

Named established theorems, such as CN, BD, HD, Harish--Chandra formal
degree, Tamagawa and exact global Artin reciprocity, are distinguished
from these bespoke suppliers.  They can be accepted as explicit citation
leaves in a readable proof.  Expanding all classical background is a
different goal from closing the new argument's dependency chains.
An interface merely bearing the word ``literature'' is not such a
citation leaf unless its full mathematical statement has been matched
to an established theorem.

The machine-readable companion records the gates, source anchors,
quantifier obligations and binding keys.  It reports no completion
percentage and does not infer a global proof from a count of local claims.
Future work should retire a leaf only with its exact supplier or a fully
proved replacement, then propagate that closure through the gates above.
